%======================================================================
\chapter{Conclusion}
%======================================================================

%A conclusion, summarizing our findings.

% TODO: TALK ABOUT THE RESULTS FROM THE THRESHOLD CHARTS

This work investigated tumor growth using a mathematical model, incorporating the dynamics of \gls{qcc}s through an extension. We employed a \gls{dde} model to track key cell populations and analyze their response to cytotoxic drugs. The stability analysis focused on the tumor-free steady state, a critical indicator of treatment effectiveness. We further explored the basin of attraction to identify treatment regimens that effectively suppress tumor growth.

The inclusion of \gls{qcc}s significantly impacts the model's behaviour, highlighting the importance of considering this subpopulation in treatment strategies. The model suggests that therapies targeting both proliferating and quiescent cancer cells are necessary for long-term tumor control.

This framework offers potential for further development:

\begin{enumerate}
    \item \textbf{\textit{Incorporating Additional Complexities:}} Future work could explore additional biological complexities such as spatial tumor heterogeneity or the immune system's role in controlling \gls{qcc}s.
    \item \textbf{\textit{Patient-Specific Modeling:}} The model could be parameterized using patient-specific data to guide personalized treatment plans.
\end{enumerate}

Overall, this study demonstrates the value of mathematical modeling in understanding tumor growth dynamics. By incorporating biological complexities like \gls{qcc}s, the model provides a more comprehensive framework for evaluating treatment strategies and paving the way for potential advancements in cancer therapy.
